\documentclass[italian,12pt]{article} %tipo di documento

\def\Title{Verbale interno}
\def\Subtitle{Riunione interna settimanale}
\def\Author{7Last}
\def\Date{2024-05-22}
\def\Version{v1.0}

%pacchetti extra da scaricare dblfloatfix, fancyhdr
\usepackage[left=2cm, right=2cm, bottom=3cm, top=3cm]{geometry}
\usepackage{fancyhdr}%creazione header-footer
\usepackage{graphicx} %serve per inserire immagini
\graphicspath{ {../../../logo/} }
%\usepackage{dblfloatfix} %serve per posizionare gli elementi dove si vuole
\usepackage[hidelinks]{hyperref} %serve per i link
\usepackage{tikz}
\usepackage{tgadventor} % font
\usepackage[useregional=numeric,showseconds=true,showzone=false]{datetime2}
\usepackage{caption}

\usepackage{hyperref}
\usepackage{tocloft}
\usepackage{titlesec}
\usepackage{color}
\usepackage{ulem}
\usepackage{pgfplots}
\usepackage{pgf-pie}
\usepackage[italian]{babel}
\usepackage{comment}
\usepackage{tabularx}
\usepackage{longtable}
\usepackage{float}
\usepackage{amsmath}
%--------creazione di sub_sub_sub_section-------%

\titleclass{\subsubsubsection}{straight}[\subsection]

\newcounter{subsubsubsection}[subsubsection]
\renewcommand\thesubsubsubsection{\thesubsubsection.\arabic{subsubsubsection}}
\renewcommand\theparagraph{\thesubsubsubsection.\arabic{paragraph}} % optional; useful if paragraphs are to be numbered

\titleformat{\subsubsubsection}
  {\normalfont\normalsize\bfseries}{\thesubsubsubsection}{1em}{}
\titlespacing*{\subsubsubsection}
{0pt}{3.25ex plus 1ex minus .2ex}{1.5ex plus .2ex}

\makeatletter
\renewcommand\paragraph{\@startsection{paragraph}{5}{\z@}%
  {3.25ex \@plus1ex \@minus.2ex}%
  {-1em}%
  {\normalfont\normalsize\bfseries}}
\renewcommand\subparagraph{\@startsection{subparagraph}{6}{\parindent}%
  {3.25ex \@plus1ex \@minus .2ex}%
  {-1em}%
  {\normalfont\normalsize\bfseries}}
\def\toclevel@subsubsubsection{4}
\def\toclevel@paragraph{5}
\def\toclevel@paragraph{6}
\def\l@subsubsubsection{\@dottedtocline{4}{7em}{4em}}
\def\l@paragraph{\@dottedtocline{5}{10em}{5em}}
\def\l@subparagraph{\@dottedtocline{6}{14em}{6em}}
\makeatother

\setcounter{secnumdepth}{4}
\setcounter{tocdepth}{4}
%----------------------------------------------------------------%

\linespread{1.2}

\captionsetup[table]{name=Tabella}
\geometry{headsep=1.5cm}

\renewcommand{\contentsname}{Indice}%imposto il nome dell'indice
\renewcommand\familydefault{\sfdefault}

\renewcommand{\listtablename}{Indice delle tabelle}%imposto il nome della lista tabelle
\renewcommand\familydefault{\sfdefault}

\renewcommand{\listfigurename}{Indice delle immagini}%imposto il nome della lista immagini
\renewcommand\familydefault{\sfdefault}

%-------------------INIZIO DOCUMENTO--------------
\begin{document}
\newgeometry{left=2cm,right=2cm,bottom=2.1cm,top=2.1cm}
\begin{titlepage}
    \vspace*{.5cm}

    \vspace{2cm}
    {
        \centering
        {\bfseries\huge \Title\par}
        \bigbreak
        {\bfseries\large \Author\par}
        \bigbreak
        {\Version\par}
        \vfill

        \begin{tikzpicture}[remember picture,overlay]

            \fill[blue!20!black] (current page.south west) -- ++(10cm,0) -- ++(-10cm,15cm) -- cycle;
            \fill[orange] (current page.south east) -- ++(-18cm,0) -- ++(21.6cm,22cm) -- cycle;

            \clip (0,-3cm) circle (2.5cm) node (current page.center) {\includegraphics[width=5cm]{logo.jpg}};
        \end{tikzpicture}

    }

    \vfill

\end{titlepage}

\restoregeometry





















\newpage
%---------------header------------------%
\pagestyle{fancy}
\fancyhead{} % pulizia degli header
\lhead{%
\begin{tikzpicture}
    \clip (0,0) circle (0.5cm);
    \node at (0,0) {\includegraphics[width=1cm]{logo}};
\end{tikzpicture}%
}
\chead{\vspace{\fill}\Title\vspace{\fill}}
\rhead{\vspace{\fill}\Version\vspace{\fill}}
%quad è una spaziatura
%---------------------------------------%
%-----------tabella revisioni-----------%
\captionsetup[table]{list=no}
\begin{table}[!h]
	\begin{center}
		Versioni\\
		\vspace{0.5cm}
		\begin{tabular}{ l l l l l }
			\hline                                                                                       \\[-2ex]
			Ver. & Data       & Autore             & Verificatore  & Descrizione                         \\
			\\[-2ex] \hline \\[-1.5ex]
			0.4  & 2024-03-16 & Elena Ferro        & Matteo Tiozzo & Riordinamento obbiettivi di qualità \\
			0.3  & 2024-04-07 & Valerio Occhinegro & Matteo Tiozzo & Stesura documento                   \\
			0.2  & 2024-03-29 & Matteo Tiozzo      & Elena Ferro   & Modificato tabella versioni         \\
			0.1  & 2024-03-28 & Valerio Occhinegro & Matteo Tiozzo & Prima redazione                     \\
			\\[-1.5ex] \hline
		\end{tabular}
	\end{center}
\end{table}
\captionsetup[table]{list=yes}
%---------------------------------------%
\newpage
\tableofcontents
\listoftables
\listoffigures
\newpage
\section{Introduzione}
\subsection{Obiettivo del documento}
Il presente documento ha lo scopo di definire le strategie di verifica e validazione utilizzate per assicurare il corretto funzionamento dello strumento sviluppato e delle
attività che lo accompagnano.  Sarà sottoposto a revisioni continue, così da prevedere situazioni precedentemente non occorse e da seguire l'evoluzione del progetto.
\subsection{Glossario}
Il \href{https://7last.github.io/docs/rtb/documentazione-interna/glossario#glossario}{glossario\textsubscript{G}} è uno strumento utilizzato per risolvere eventuali dubbi riguardanti 
alcuni termini specifici utilizzati nella redazione del documento.
Esso conterrà la definizione dei termini evidenziati e sarà consultabile al seguente \href{https://7last.github.io/docs/rtb/documentazione-interna/glossario}{link}. I termini presenti in tale documento saranno evidenziati da una 'G' a pedice.
\subsection{Riferimenti}
\subsubsection{Riferimenti normativi}
\begin{itemize}
    \item \href{https://7last.github.io/docs/rtb/documentazione-interna/glossario#norme-di-progetto}{Norme di progetto\textsubscript{G}} (aggiungere versione e/o link al documento);
    \item Regolamento del progetto:\\
		  \url{https://www.math.unipd.it/~tullio/IS-1/2023/Dispense/PD2.pdf}.
\end{itemize}
\subsubsection{Riferimenti informativi}
\begin{itemize}
    \item \href{https://7last.github.io/docs/rtb/documentazione-interna/glossario#capitolato}{Capitolato\textsubscript{G}} d'appalto C6: \href{https://7last.github.io/docs/rtb/documentazione-interna/glossario#synccity}{SyncCity\textsubscript{G}} – A \href{https://7last.github.io/docs/rtb/documentazione-interna/glossario#smart-city}{smart city\textsubscript{G}} monitoring platform\\
    \url{https://www.math.unipd.it/~tullio/IS-1/2023/Progetto/C6.pdf};
    \item \href{https://it.wikipedia.org/wiki/ISO/IEC_9126}{Standard ISO/IEC 9126};
    \item \href{https://iso25000.com/index.php/en/iso-25000-standards/iso-25010}{Standard ISO/IEC 25010};
    \item \href{ https://en.wikipedia.org/wiki/ISO/IEC_12207}{Standard ISO/IEC 12207:1995};
    \item \href{URL}{\textit{Verbali esterni}};
    \item \href{URL}{\textit{Verbali interni}};
    \item \href{URL}{\href{https://7last.github.io/docs/rtb/documentazione-interna/glossario#analisi-dei-requisiti}{\textit{Analisi dei requisiti}\textsubscript{G}}};
    \item AGGIUNGERE LINK
\end{itemize}
\include{metriche_qualità.tex}
\section{Metodologie di Testing}
La fase di testing è un'attività fondamentale per garantire la qualità del prodotto software. Permette di verificare che il software sia conforme ai requisiti e alle specifiche richieste e di individuare tempestivamente eventuali bug o problemi di funzionamento, così da poterli correggere prima del rilascio del prodotto; garantisce inoltre che gli stessi errori non si ripetano in futuro. \\
In questa sezione verranno descritte le metodologie di testing adottate per garantire il rispetto dei vincoli individuati nella sezione \textit{Requisiti} del documento \href{https://7last.github.io/docs/rtb/documentazione-esterna/analisi-dei-requisiti}{\href{https://7last.github.io/docs/rtb/documentazione-interna/glossario\#analisi-dei-requisiti}{\textit{Analisi dei Requisiti}\textsubscript{G}}}. \\ % TODO aggiungere versione documento
Nelle successive sottosezioni verranno descritte le tipologie di test effettuati con l'indicazione del codice del test, una breve descrizione di ciò che viene verificato e lo stato di superamento del test, espresso nel seguente modo: \\
\begin{itemize}
	\item \textbf{S}: test superato;
	\item \textbf{NS}: test non superato;
	\item \textbf{NI}: test non implementato.
\end{itemize}

\subsection{Test di Unità}
I test di unità sono test che verificano il corretto funzionamento delle singole unità di codice, ovvero le più piccole parti di un programma. Questi test vengono effettuati per verificare che ogni unità funzioni correttamente e che sia in grado di eseguire le operazioni richieste. \\
\begin{longtable}{|>{\raggedright\arraybackslash}m{0.1\textwidth}|>{\raggedright\arraybackslash}m{0.6\textwidth}|>{\raggedright\arraybackslash}m{0.1\textwidth}|}
	\hline
	\textbf{Codice} & \textbf{Descrizione}                                                                         & \textbf{Stato} \\
	\hline
	\endfirsthead
	\hline
	\textbf{Codice} & \textbf{Descrizione}                                                                         & \textbf{Stato} \\
	\endhead
	\textbf{1T-U}   & Verificare che la classe \textit{temperature} venga creata correttamente.                    & NI             \\
	\hline
	\textbf{2T-U}   & Verificare che il livello di salute di una città si modifichi ad ogni nuova misurazione.     & NI             \\
	\hline
	\textbf{3T-U}   & Verificare che il metodo "GENERA\_MISURA" di una classe generi correttamente la misurazione. & NI             \\
	\hline
	\caption{Test di Unità}
\end{longtable}

\subsection{Test di Integrazione}
I test di integrazione sono test che verificano il corretto funzionamento delle interfacce tra le varie unità di codice. Questi test vengono effettuati per verificare che le varie unità di codice e i vari moduli interagiscano correttamente tra di loro e che siano in grado di comunicare e scambiarsi i dati necessari. \\
\begin{longtable}{|>{\raggedright\arraybackslash}m{0.1\textwidth}|>{\raggedright\arraybackslash}m{0.6\textwidth}|>{\raggedright\arraybackslash}m{0.1\textwidth}|}
	\hline
	\textbf{Codice} & \textbf{Descrizione}                                                                                    & \textbf{Stato} \\
	\hline
	\endfirsthead
	\hline
	\textbf{Codice} & \textbf{Descrizione}                                                                                    & \textbf{Stato} \\
	\endhead
	\textbf{1T-I}   & Verificare che i dati generati dal sensore di temperatura siano memorizzati correttamente nel database. & NI             \\
	\hline
	\textbf{2T-I}   & Verificare che i dati generati dal sensore di traffico siano memorizzati correttamente nel database.    & NI             \\
	\hline
	\textbf{3T-I}   & Verificare che i dati della temperatura media siano memorizzati correttamente nel database.             & NI             \\
	\hline
	\textbf{4T-I}   & Verificare che i dati della velocità media dei veicoli siano memorizzati correttamente nel database.    & NI             \\
	\hline
	\textbf{5T-I}   & Verificare che i dati dei veicoli transitati siano memorizzati correttamente.                           & NI             \\
	\hline
	\caption{Test di Integrazione}
\end{longtable}

\subsection{Test di Sistema}
I test di sistema sono finalizzati alla verifica del soddisfacimento dei requisiti richiesti ed evidenziati nel documento \href{https://7last.github.io/docs/rtb/documentazione-esterna/analisi-dei-requisiti}{\href{https://7last.github.io/docs/rtb/documentazione-interna/glossario\#analisi-dei-requisiti}{\textit{Analisi dei Requisiti}\textsubscript{G}}}. Questi test vengono effettuati sul sistema nel suo complesso, per verificare che il software funzioni correttamente e che sia in grado di eseguire le operazioni richieste. \\
\begin{longtable}{|>{\raggedright\arraybackslash}m{0.1\textwidth}|>{\raggedright\arraybackslash}m{0.6\textwidth}|>{\raggedright\arraybackslash}m{0.1\textwidth}|}
	\hline
	\textbf{Codice} & \textbf{Descrizione}                                                                                                                                                               & \textbf{Stato} \\
	\hline
	\endfirsthead
	\hline
	\textbf{Codice} & \textbf{Descrizione}                                                                                                                                                               & \textbf{Stato} \\
	\endhead
	\textbf{1T-S}   & Verificare che l'accesso al sistema non richieda alcuna procedura di login e che sia direttamente accessibile dall'utente.                                                         & NI             \\
	\hline
	\textbf{2T-S}   & Verificare che il prodotto non abbia alcuna sezione o funzionalità di amministrazione o gestione riservata.                                                                        & NI             \\
	\hline
	\textbf{3T-S}   & Verificare che i sensori integrati producano una misurazione coerente con il tipo di sensore simulato.                                                                             & NI             \\
	\hline
	\textbf{4T-S}   & Verificare che ogni misurazione inviata dal simulatore contenga l’identificativo del sensore, la misurazione d'interesse e il timestamp.                                           & NI             \\
	\hline
	\textbf{5T-S}   & Verificare che il sistema sia in grado di ricevere e memorizzare correttamente le misurazioni inviate dai sensori.                                                                 & NI             \\
	\hline
	\textbf{6T-S}   & Verificare che il sistema sia in grado di simulare almeno un sensore per rilevare la temperatura.                                                                                  & NI             \\
	\hline
	\textbf{7T-S}   & Verificare che ogni dato generato dai simulatori dei sensori sia strettamente correlato al dato successivo, garantendo una transizione realistica tra le misurazioni.              & NI             \\
	\hline
	\textbf{8T-S}   & Verificare che il prodotto di visualizzazione supporti la rappresentazione di dati provenienti da diversi tipi di sensori, permettendo una simulazione quanto più possibile reale. & NI             \\
	\hline
	\textbf{9T-S}   & Verificare che l'utente possa vedere una dashboard completo delle condizioni attuali della città tramite l'uso di appositi widget rappresentanti le misurazioni dei sensori.       & NI             \\
	\hline
	\textbf{10T-S}  & Verificare la facilità di comprensione e l'intuitività dell'interfaccia grafica, garantendo un'esperienza utente piacevole e soddisfacente.                                        & NI             \\
	\hline
	\textbf{11T-S}  & Verificare che il sistema sia in grado di inviare notifiche all'utente in caso di superamento di soglie di temperatura predefinite.                                                & NI             \\
	\hline
	\textbf{12T-S}  & Verificare che l’utente possa vedere le misurazioni all’interno dei widget dedicati alla rappresentazione delle rilevazioni dei sensori in un formato testuale.                    & NI             \\
	\hline
	\textbf{13T-S}  & Verificare che la dashboard si aggiorni quasi istantaneamente per riflettere i dati provenienti dai sensori entro un massimo di 15 secondi.                                        & NI             \\
	\hline
	\textbf{14T-S}  & Verificare che ogni widget che visualizza le misurazioni includa informazioni sull’identificativo dei sensori che hanno contribuito a quelle misurazioni.                          & NI             \\
	\hline
	\textbf{15T-S}  & Verificare che la dashboard contenga almeno un widget dedicato alle misurazioni dei sensori di temperatura.                                                                        & NI             \\
	\hline
	\textbf{16T-S}  & Verificare che la dashboard contenga almeno un widget dedicato alle misurazioni dei sensori di traffico.                                                                           & NI             \\
	\hline
	\textbf{17T-S}  & Verificare che i sensori presenti sulla mappa siano distinguibili in modo chiaro ed inequivocabile, permettendo il riconoscimento della loro tipologia.                            & NI             \\
	\hline
	\textbf{18T-S}  & Verificare che l’utente possa filtrare la visualizzazione delle misurazioni di una specifica tipologia di sensori.                                                                 & NI             \\
	\hline
	\textbf{19T-S}  & Verificare che il sistema verifichi la validità del sensore inserito dall’utente.                                                                                                  & NI             \\
	\hline
	\textbf{20T-S}  & Verificare che, in caso di inserimento di un sensore non valido, il sistema generi una notifica di errore.                                                                         & NI             \\
	\hline
	\textbf{21T-S}  & Verificare che la notifica di errore relativa all’inserimento di un sensore non valido richieda all’utente di inserire nuovamente il sensore.                                      & NI             \\
	\hline
	\textbf{22T-S}  & Verificare che la notifica di errore relativa all’inserimento di un sensore non valido sia chiara e informativa, indicando il motivo specifico dell’invalidità del sensore.        & NI             \\
	\hline
	\textbf{23T-S}  & Verificare che l’utente possa visualizzare la lista delle misurazioni rilevanti.                                                                                                   & NI             \\
	\hline
	\textbf{24T-S}  & Verificare che ogni misurazione nella lista dei più importanti fornisca correttamente l’identificativo del sensore.                                                                & NI             \\
	\hline
	\textbf{25T-S}  & Verificare che ogni misurazione nella lista dei più importanti fornisca correttamente la tipologia del sensore.                                                                    & NI             \\
	\hline
	\textbf{26T-S}  & Verificare che ogni misurazione nella lista dei più importanti fornisca correttamente il valore della misurazione.                                                                 & NI             \\
	\hline
	\textbf{27T-S}  & Verificare che ogni misurazione nella lista dei più importanti fornisca correttamente il timestamp della misurazione.                                                              & NI             \\
	\hline
	\textbf{28T-S}  & Verificare che l’utente possa visualizzare la lista delle misurazioni rilevanti in un formato tabellare.                                                                           & NI             \\ %da controllare
	\hline
	\textbf{29T-S}  & Verificare che l’utente possa rimuovere una misurazione specifica dalla lista delle misurazioni rilevanti.                                                                         & NI             \\
	\hline
	\textbf{30T-S}  & Verificare che l’utente possa rimuovere tutte le misurazioni dalla lista delle misurazioni rilevanti.                                                                              & NI             \\
	\hline
	\textbf{31T-S}  & Verificare che l’utente riceva notifiche quando i sensori superano predeterminate soglie di sicurezza.                                                                             & NI             \\
	\hline
	\textbf{32T-S}  & Verificare che l’utente possa visualizzare le notifiche relative ai sensori che superano le soglie di sicurezza.                                                                   & NI             \\
	\hline
	\textbf{33T-S}  & Verificare che l’utente possa visualizzare correttamente le informazioni richieste per i sensori.                                                                                  & NI             \\
	\hline
	\textbf{34T-S}  & Verificare che l’utente possa visualizzare correttamente la posizione in coordinate dei sensori.                                                                                   & NI             \\
	\hline
	\textbf{35T-S}  & Verificare che l’utente possa visualizzare correttamente la data di installazione del sensore .                                                                                    & NI             \\
	\hline
	\textbf{36T-S}  & Verificare che l’utente possa visualizzare correttamente la data di ultima manutenzione del sensore.                                                                               & NI             \\
	\hline
	\textbf{37T-S}  & Verificare che l’utente possa visualizzare correttamente l’unità di misura associata al sensore.                                                                                   & NI             \\
	\hline
	\caption{Test di Sistema}
\end{longtable}

\subsection{Test di Regressione}
I test di regressione sono test che vengono effettuati per verificare che le modifiche apportate al software non abbiano introdotto nuovi errori o problemi di funzionamento e che il software continui a funzionare correttamente anche dopo le modifiche fatte. \\
\begin{longtable}{|>{\raggedright\arraybackslash}m{0.1\textwidth}|>{\raggedright\arraybackslash}m{0.6\textwidth}|>{\raggedright\arraybackslash}m{0.1\textwidth}|}
	\hline
	\textbf{Tipologia di test}    & \textbf{Codice}                           & \textbf{Stato} \\
	\hline
	\endfirsthead
	\hline
	\textbf{Tipologia di test}    & \textbf{Codice}                           & \textbf{Stato} \\
	\endhead
	\textbf{Test di unità}        & Codice del test che andremo ad effettuare & NI             \\
	\hline
	\textbf{Test di sistema}      & Codice del test che andremo ad effettuare & NI             \\
	\hline
	\textbf{Test di integrazione} & Codice del test che andremo ad effettuare & NI             \\
	\hline
	\caption{Test di Regressione}
\end{longtable}
\newpage
\subsection{Test di Accettazione}
I test di accettazione sono test che vengono effettuati per verificare che il software soddisfi i requisiti richiesti. Consentono di ultimare il processo di validazione del prodotto finale. Questi test verranno eseguiti sia dal gruppo di sviluppo \textit{7Last} che dall'azienda \href{https://7last.github.io/docs/rtb/documentazione-interna/glossario\#proponente}{proponente\textsubscript{G}} \textit{Sync Lab S.r.l.}. \\
\begin{longtable}{|>{\raggedright\arraybackslash}m{0.1\textwidth}|>{\raggedright\arraybackslash}m{0.6\textwidth}|>{\raggedright\arraybackslash}m{0.1\textwidth}|}
	\hline
	\textbf{Codice} & \textbf{Descrizione}                                                                                                                                                  & \textbf{Stato} \\
	\hline
	\endfirsthead
	\hline
	\textbf{Codice} & \textbf{Descrizione}                                                                                                                                                  & \textbf{Stato} \\
	\endhead
	\textbf{1T-A}   & Verificare che tutti i widget relativi alle diverse tipologie di sensori siano visibili sulla dashboard.                                                              & NI             \\
	\hline
	\textbf{2T-A}   & Verificare che la mappa dei sensori si carichi correttamente e permetta interazioni fluide.                                                                           & NI             \\
	\hline
	\textbf{3T-A}   & Verificare che le informazioni di un sensore specifico siano visualizzate correttamente quando selezionate dalla dashboard.                                           & NI             \\
	\hline
	\textbf{4T-A}   & Verifica della gestione corretta degli errori nel caso in cui i dati dei sensori non siano disponibili o siano incompleti all’interno della visualizzazione testuale. & NI             \\
	\hline
	\textbf{5T-A}   & Verifica della corretta visualizzazione delle misurazioni effettuate nel tempo dai sensori.                                                                           & NI             \\
	\hline
	\textbf{6T-A}   & Verificare l’accuratezza e la completezza delle opzioni di interazione offerte dall’interfaccia del widget per esaminare i dati di temperatura.                       & NI             \\
	\hline
	\textbf{7T-A}   & Verificare ci sia l’opportunità di visualizzare correttamente il widget contenente le misurazioni dei sensori di temperatura.                                         & NI             \\
	\hline
	\textbf{8T-A}   & Verificare si possa filtrare correttamente la visualizzazione delle misurazioni in base ad una specifica selezione di sensori.                                        & NI             \\
	\hline
	\textbf{9T-A}   & Verificare che si possa rimuovere correttamente i filtri attivi per visualizzazione delle misurazioni dei sensori.                                                    & NI             \\
	\hline
	\textbf{10T-A}  & Verificare che si riceva correttamente una notifica in caso di superamento delle soglie impostate per le misurazioni.                                                 & NI             \\
	\hline
	\textbf{11T-A}  & Verificare che si possa visualizzare correttamente la lista delle misurazioni rilevanti.                                                                              & NI             \\
	\hline
	\textbf{12T-A}  & Verificare che si possa rimuovere correttamente una misurazione specifica dalla lista delle misurazioni rilevanti.                                                    & NI             \\
	\hline
	\textbf{13T-A}  & Verificare che si possa rimuovere correttamente tutte le misurazioni dalla lista delle misurazioni rilevanti.                                                         & NI             \\
	\hline
	\textbf{14T-A}  & Verificare che si possa inserire correttamente una misurazione nella lista delle misurazioni rilevanti.                                                               & NI             \\
	\hline
	\caption{Test di Accettazione}
\end{longtable}

\definecolor{opt}{HTML}{51B824}
\definecolor{amm}{HTML}{CC0000}

\section{Cruscotto di valutazione della qualità}
\subsection{Qualità del processo di Analisi dei requisiti}
\subsubsection{1M-CRO - Copertura dei requisiti obbligatori}
%--------- GRAFICO -----------%
\begin{figure*}[!h]
    \centering
    \begin{tikzpicture}
        \begin{axis}[
            width  = 0.85*\textwidth,
            height = 8cm,
            ymajorgrids = true,
            symbolic x coords={Sprint 1, Sprint 2, Sprint 3},
            xtick = data,
            ytick = {0, 25, 50, 75, 100},
            ymin=0, ymax=100,
            axis lines*=left,
            legend cell align=left,
            legend style={
                at={(0.5,1.15)},
                anchor=south,
                column sep=0.1ex,
                legend columns=3
            },      
            xticklabel style={rotate=45, anchor=north east, yshift=0ex, xshift=0ex}
            ]
            \addplot[color=opt, style={dashed, line width=3pt}, mark=none] coordinates {(Sprint 1, 100) (Sprint 2, 100) (Sprint 3, 100)};
            \addplot[color=blue, style={line width=1pt}, mark=none] coordinates {(Sprint 1, 25) (Sprint 2, 48) (Sprint 3, 52)};
            \legend{Valore ottimo ed ammissibile, Valore effettivo}
        \end{axis}
    \end{tikzpicture}
    %44 requisiti obbligatori, 11 implementati nel primo sprint, 21 nel secondo e 24 nel terzo
    \caption{Percentuale di copertura dei requisiti obbligatori}
    \label{fig:1}
\end{figure*}
%--------- FINE GRAFICO -----------%
%\textbf{RTB} \\
% TODO considerazioni finali per RTB \\
%\textbf{PB} \\
% TODO considerazioni finali per PB

\newpage
\subsubsection{2M-CRD - Copertura dei requisiti desiderabili}
%--------- GRAFICO -----------%
\begin{figure*}[!h]
    \centering
    \begin{tikzpicture}
        \begin{axis}[
            width  = 0.85*\textwidth,
            height = 8cm,
            ymajorgrids = true,
            symbolic x coords={Sprint 1, Sprint 2, Sprint 3},
            xtick = data,
            ytick = {0, 25, 50, 75, 100},
            ymin=0, ymax=100,
            axis lines*=left,
            legend cell align=left,
            legend style={
                at={(0.5,1.15)},
                anchor=south,
                column sep=0.1ex,
                legend columns=3
            },      
            xticklabel style={rotate=45, anchor=north east, yshift=0ex, xshift=0ex}
            ]
            \addplot[color=opt, style={dashed, line width=3pt}, mark=none] coordinates {(Sprint 1, 100) (Sprint 2, 100) (Sprint 3, 100)};
            \addplot[color=amm, style={dashed, line width=3pt}, mark=none] coordinates {(Sprint 1, 35) (Sprint 2, 35) (Sprint 3, 35)};
            \addplot[color=blue, style={line width=1pt}, mark=none] coordinates {(Sprint 1, 50) (Sprint 2, 100) (Sprint 3, 100)};
            \legend{Valore ottimo, Valore ammissibile, Valore effettivo}
        \end{axis}
    \end{tikzpicture}
    \caption{Percentuale di copertura dei requisiti desiderabili}
    \label{fig:2} 
\end{figure*}
%--------- FINE GRAFICO -----------%
%\textbf{RTB} \\
% TODO considerazioni finali per RTB \\
%\textbf{PB} \\
% TODO considerazioni finali per PB

\newpage
\subsubsection{3M-CROP - Copertura dei requisiti opzionali}
%--------- GRAFICO -----------%
\begin{figure*}[!h]
    \centering
    \begin{tikzpicture}
        \begin{axis}[
            width  = 0.85*\textwidth,
            height = 8cm,
            ymajorgrids = true,
            symbolic x coords={Sprint 1, Sprint 2, Sprint 3},
            xtick = data,
            ytick = {0, 25, 50, 75, 100},
            ymin=0, ymax=100,
            axis lines*=left,
            legend cell align=left,
            legend style={
                at={(0.5,1.15)},
                anchor=south,
                column sep=0.1ex,
                legend columns=3
            },      
            xticklabel style={rotate=45, anchor=north east, yshift=0ex, xshift=0ex}
            ]
            \addplot[color=opt, style={dashed, line width=3pt}, mark=none] coordinates {(Sprint 1, 100) (Sprint 2, 100) (Sprint 3, 100)};
            \addplot[color=amm, style={dashed, line width=3pt}, mark=none] coordinates {(Sprint 1, 0) (Sprint 2, 0) (Sprint 3, 0)};
            \addplot[color=blue, style={line width=1pt}, mark=none] coordinates {(Sprint 1, 0) (Sprint 2, 0) (Sprint 3, 0)};
            \legend{Valore ottimo, Valore ammissibile, Valore effettivo}
        \end{axis}
    \end{tikzpicture}
    \caption{Percentuale di copertura dei requisiti opzionali} %non abbiamo implementato nessun requisito opzionale
    \label{fig:3}
\end{figure*}
%--------- FINE GRAFICO -----------%
%\textbf{RTB} \\
% TODO considerazioni finali per RTB \\
%\textbf{PB} \\
% TODO considerazioni finali per PB

\newpage
\subsection{Qualità del processo di Fornitura}
\subsubsection{9M-EV - Earned Value (EV) e 10M-PV - Planned Value (PV)}
%--------- GRAFICO -----------%
\begin{figure*}[!h]
    \centering
    \begin{tikzpicture}
        %da chiedere se è giusta questa interpretazione
        \begin{axis}[
            width  = 0.85*\textwidth,
            height = 8cm,
            ymajorgrids = true,
            symbolic x coords={Sprint 1, Sprint 2, Sprint 3},
            xtick = data,
            ymin=0,
            axis lines*=left,
            legend cell align=left,
            legend style={
                at={(0.5,1.15)},
                anchor=south,
                column sep=0.1ex,
                legend columns=3
            },      
            xticklabel style={rotate=45, anchor=north east, yshift=0ex, xshift=0ex},
            scaled y ticks = false,
            yticklabel style={/pgf/number format/fixed}
            ]
            \addplot[color=opt, style={dashed, line width=3pt}, mark=none] coordinates {(Sprint 1, 12670) (Sprint 2, 11730) (Sprint 3, 12660)}; % TODO mettere valori corretti EAC
            \addplot[color=amm, style={dashed, line width=3pt}, mark=none] coordinates {(Sprint 1, 1090) (Sprint 2, 2220) (Sprint 3, 2925)}; % TODO mettere valori corretti PV
            \addplot[color=blue, style={line width=1pt}, mark=none] coordinates {(Sprint 1, 977.5) (Sprint 2, 2032.5) (Sprint 3, 2792.5)}; % TODO mettere valori corretti EV
            \legend{EAC, PV, EV}
        \end{axis}
    \end{tikzpicture}
    \caption{Proiezione del PV e dell'EV}
    \label{fig:4}
\end{figure*}
%--------- FINE GRAFICO -----------%
%\textbf{RTB} \\
% TODO considerazioni finali per RTB \\
%\textbf{PB} \\
% TODO considerazioni finali per PB

\newpage
\subsubsection{11M-AC - Actual Cost (AC) e 14M-ETC - Estimate to Complete (ETC)}
% da controllare (sistemare ETC)
%--------- GRAFICO -----------%
\begin{figure*}[!h]
    \centering
    \begin{tikzpicture}
        \begin{axis}[
            width  = 0.85*\textwidth,
            height = 8cm,
            ymajorgrids = true,
            symbolic x coords={Sprint 1, Sprint 2, Sprint 3},
            xtick = data,
            ymin=0,
            axis lines*=left,
            legend cell align=left,
            legend style={
                at={(0.5,1.15)},
                anchor=south,
                column sep=0.1ex,
                legend columns=3
            },      
            xticklabel style={rotate=45, anchor=north east, yshift=0ex, xshift=0ex},
            scaled y ticks = false,
            yticklabel style={/pgf/number format/fixed}
            ]
            \addplot[color=opt, style={dashed, line width=3pt}, mark=none] coordinates {(Sprint 1, 12670) (Sprint 2, 11730) (Sprint 3, 12660)}; 
            \addplot[color=amm, style={dashed, line width=3pt}, mark=none] coordinates {(Sprint 1, 11695) (Sprint 2, 10670) (Sprint 3, 9670)};
            \addplot[color=blue, style={line width=1pt}, mark=none] coordinates {(Sprint 1, 977.5) (Sprint 2, 2032.5) (Sprint 3, 2792.5)};
            \legend{EAC, ETC, AC}
        \end{axis}
    \end{tikzpicture}
    \caption{Proiezione dell'AC e dell'ETC}
    \label{fig:5}
\end{figure*}
%--------- FINE GRAFICO -----------%
%\textbf{RTB} \\
% TODO considerazioni finali per RTB \\
%\textbf{PB} \\
% TODO considerazioni finali per PB

\newpage
\subsubsection{12M-CV - Cost Variance (CV) e 33M-SV - Schedule Variance (SV)}
%--------- GRAFICO -----------%
\begin{figure*}[!h]
    \centering
    \begin{tikzpicture}
        \begin{axis}[
            width  = 0.85*\textwidth,
            height = 8cm,
            ymajorgrids = true,
            symbolic x coords={Sprint 1, Sprint 2, Sprint 3},
            xtick = data,
            ymin=-10, ymax=10,
            axis lines*=left,
            legend cell align=left,
            legend style={
                at={(0.5,1.15)},
                anchor=south,
                column sep=0.1ex,
                legend columns=3
            },      
            xticklabel style={rotate=45, anchor=north east, yshift=0ex, xshift=0ex}
            ]
            \addplot[color=opt, style={dashed, line width=3pt}, mark=none] coordinates {(Sprint 1, 0) (Sprint 2, 0) (Sprint 3, 0)};
            \addplot[color=amm, style={dashed, line width=3pt}, mark=none] coordinates {(Sprint 1, -7.5) (Sprint 2, -7.5) (Sprint 3, -7.5)};
            \addplot[color=blue, style={line width=1pt}, mark=none] coordinates {(Sprint 1, 4) (Sprint 2, 1) (Sprint 3, -3)}; % TODO mettere valori corretti SV terzo periodo
            \addplot[color=yellow, style={line width=1pt}, mark=none] coordinates {(Sprint 1, 9) (Sprint 2, 7) (Sprint 3, -8)}; % TODO mettere valori corretti CV terzo periodo
            \legend{Valore ottimo, Valore ammissibile, SV, CV}
        \end{axis}
    \end{tikzpicture}
    \caption{Andamento percentuale di SV e CV}
    \label{fig:6}
\end{figure*}
%--------- FINE GRAFICO -----------%
%\textbf{RTB} \\
% TODO considerazioni finali per RTB \\
%\textbf{PB} \\
% TODO considerazioni finali per PB

\newpage
\subsubsection{13M-EAC - Estimated at Completion (EAC)}
%--------- GRAFICO -----------%
\begin{figure*}[!h]
    \centering
    \begin{tikzpicture}
        \begin{axis}[
            width  = 0.85*\textwidth,
            height = 8cm,
            ymajorgrids = true,
            symbolic x coords={Sprint 1, Sprint 2, Sprint 3},
            xtick = data,
            ytick = {10000, 11000, 12000, 13000, 14000, 15000},
            ymin=10000, ymax=15000,
            axis lines*=left,
            legend cell align=left,
            legend style={
                at={(0.5,1.15)},
                anchor=south,
                column sep=0.1ex,
                legend columns=3
            },      
            xticklabel style={rotate=45, anchor=north east, yshift=0ex, xshift=0ex},
            scaled y ticks = false,
            yticklabel style={/pgf/number format/fixed}
            ]
            \addplot[color=opt, style={dashed, line width=3pt}, mark=none] coordinates {(Sprint 1, 12670) (Sprint 2, 12670) (Sprint 3, 12670)}; % TODO mettere valori corretti BAC
            \addplot[color=amm, style={dashed, line width=3pt}, mark=none] coordinates {(Sprint 1, 13177) (Sprint 2, 13177) (Sprint 3, 13177)}; % TODO mettere valori corretti BAC + 4%
            \addplot[color=amm, style={dashed, line width=3pt}, mark=none] coordinates {(Sprint 1, 12162) (Sprint 2, 12162) (Sprint 3, 12162)}; % TODO mettere valori corretti BAC - 4%
            \addplot[color=blue, style={line width=1pt}, mark=none] coordinates {(Sprint 1, 11730) (Sprint 2, 11605) (Sprint 3, 11000)}; % TODO mettere valori corretti EAC
            \legend{Valore ottimo (BAC), Valore ammissibile superiore, Valore ammissibile inferiore, Valore effettivo}
        \end{axis}
    \end{tikzpicture}
    \caption{Proiezione dell'EAC}
    \label{fig:7}
\end{figure*}
%--------- FINE GRAFICO -----------%
%\textbf{RTB} \\
% TODO considerazioni finali per RTB \\
%\textbf{PB} \\
% TODO considerazioni finali per PB

\newpage
\subsection{Qualità del processo di Documentazione}
\subsubsection{22M-IG - Indice Gulpease}
%--------- GRAFICO -----------%
\begin{figure*}[!h]
    \centering
    \begin{tikzpicture}
        \begin{axis}[
            width  = 0.85*\textwidth,
            height = 8cm,
            ymajorgrids = true,
            symbolic x coords={Sprint 1, Sprint 2, Sprint 3},
            xtick = data,
            ytick = {0, 25, 50, 75, 100},
            ymin=0, ymax=100,
            axis lines*=left,
            legend cell align=left,
            legend style={
                at={(0.5,1.15)},
                anchor=south,
                column sep=0.1ex,
                legend columns=3
            },           
            xticklabel style={rotate=45, anchor=north east, yshift=0ex, xshift=0ex}
            ]
            \addplot[color=opt, style={dashed, line width=3pt}, mark=none] coordinates {(Sprint 1, 90) (Sprint 2, 90) (Sprint 3, 90)};
            \addplot[color=amm, style={dashed, line width=3pt}, mark=none] coordinates {(Sprint 1, 60) (Sprint 2, 60) (Sprint 3, 60)};
            \addplot[color=cyan, style={line width=1pt}, mark=none] coordinates {(Sprint 1, 55) (Sprint 2, 60) (Sprint 3, 60)}; % TODO mettere valori corretti NdP
            \addplot[color=magenta, style={line width=1pt}, mark=none] coordinates {(Sprint 1, 70) (Sprint 2, 55) (Sprint 3, 55)}; % TODO mettere valori corretti PdP
            \addplot[color=lime, style={line width=1pt}, mark=none] coordinates {(Sprint 1, 80) (Sprint 2, 70) (Sprint 3, 70)}; % TODO mettere valori corretti PdQ
            \addplot[color=brown, style={line width=1pt}, mark=none] coordinates {(Sprint 1, 85) (Sprint 2, 90) (Sprint 3, 90)}; % TODO mettere valori corretti Glossario
            \addplot[color=pink, style={line width=1pt}, mark=none] coordinates {(Sprint 1, 80) (Sprint 2, 75) (Sprint 3, 75)}; % TODO mettere valori corretti AdR
            \addplot[color=violet, style={line width=1pt}, mark=none] coordinates {(Sprint 1, 0) (Sprint 2, 0) (Sprint 3, 0)}; % TODO mettere valori corretti MU
            \addplot[color=yellow, style={line width=1pt}, mark=none] coordinates {(Sprint 1, 0) (Sprint 2, 0) (Sprint 3, 0)}; % TODO mettere valori corretti ST
            \legend{Valore ottimo, Valore ammissibile, NdP, PdP, PdQ, Glossario, AdR, MU, ST}
        \end{axis}
    \end{tikzpicture}
    \caption{Andamento indice di Gulpease per ciascun documento}
    \label{fig:8}
\end{figure*}
%--------- FINE GRAFICO -----------%
%\textbf{RTB} \\
% TODO considerazioni finali per RTB \\
%\textbf{PB} \\
% TODO considerazioni finali per PB

\newpage
\subsubsection{23M-CO - Correttezza Ortografica}
%--------- GRAFICO -----------%
\begin{figure*}[!h]
    \centering
    \begin{tikzpicture}
        \begin{axis}[
            width  = 0.85*\textwidth,
            height = 8cm,
            ymajorgrids = true,
            symbolic x coords={Sprint 1, Sprint 2, Sprint 3},
            xtick = data,
            ymin=0, ymax=10,
            axis lines*=left,
            legend cell align=left,
            legend style={
                at={(0.5,1.15)},
                anchor=south,
                column sep=0.1ex,
                legend columns=3
            },           
            xticklabel style={rotate=45, anchor=north east, yshift=0ex, xshift=0ex}
            ]
            \addplot[color=opt, style={dashed, line width=3pt}, mark=none] coordinates {(Sprint 1, 0) (Sprint 2, 0) (Sprint 3, 0)};
            \addplot[color=cyan, style={line width=1pt}, mark=none] coordinates {(Sprint 1, 9) (Sprint 2, 5) (Sprint 3, 5)}; % TODO mettere valori corretti NdP
            \addplot[color=magenta, style={line width=1pt}, mark=none] coordinates {(Sprint 1, 8) (Sprint 2, 7) (Sprint 3, 4)}; % TODO mettere valori corretti PdP
            \addplot[color=lime, style={line width=1pt}, mark=none] coordinates {(Sprint 1, 9) (Sprint 2, 3) (Sprint 3, 3)}; % TODO mettere valori corretti PdQ
            \addplot[color=brown, style={line width=1pt}, mark=none] coordinates {(Sprint 1, 0) (Sprint 2, 0) (Sprint 3, 0)}; % TODO mettere valori corretti Glossario
            \addplot[color=pink, style={line width=1pt}, mark=none] coordinates {(Sprint 1, 0) (Sprint 2, 6) (Sprint 3, 3)}; % TODO mettere valori corretti AdR
            \addplot[color=violet, style={line width=1pt}, mark=none] coordinates {(Sprint 1, 0) (Sprint 2, 0) (Sprint 3, 0)}; % TODO mettere valori corretti MU
            \addplot[color=yellow, style={line width=1pt}, mark=none] coordinates {(Sprint 1, 0) (Sprint 2, 0) (Sprint 3, 0)}; % TODO mettere valori corretti ST
            \legend{Valore ottimo ed ammissibile, NdP, PdP, PdQ, Glossario, AdR, MU, ST}
        \end{axis}
    \end{tikzpicture}
    \caption{Errori ortografici per ciascun documento}
    \label{fig:9}
\end{figure*}
%--------- FINE GRAFICO -----------%
%\textbf{RTB} \\
% TODO considerazioni finali per RTB \\
%\textbf{PB} \\
% TODO considerazioni finali per PB

\newpage
\subsection{Qualità del processo di Verifica}
\subsubsection{24M-CC - Code coverage}
%Da qui in poi byteops fanno da "periodo 6"
%--------- GRAFICO -----------%
\begin{figure*}[!h]
    \centering
    \begin{tikzpicture}
        \begin{axis}[
            width  = 0.85*\textwidth,
            height = 8cm,
            ymajorgrids = true,
            symbolic x coords={Sprint 1, Sprint 2, Sprint 3},
            xtick = data,
            ytick = {0, 25, 50, 75, 100},
            ymin=0, ymax=100,
            axis lines*=left,
            legend cell align=left,
            legend style={
                at={(0.5,1.15)},
                anchor=south,
                column sep=0.1ex,
                legend columns=-1
            },           
            xticklabel style={rotate=45, anchor=north east, yshift=0ex, xshift=0ex}
            ]
            \addplot[color=opt, style={dashed, line width=3pt}, mark=none] coordinates {(Sprint 1, 100) (Sprint 2, 100) (Sprint 3, 100)};
            \addplot[color=amm, style={dashed, line width=3pt}, mark=none] coordinates {(Sprint 1, 80) (Sprint 2, 80) (Sprint 3, 80)};
            \addplot[color=blue, style={line width=1pt}, mark=none] coordinates {(Sprint 1, 0) (Sprint 2, 0) (Sprint 3, 0)};
            \legend{Valore ottimo, Valore ammissibile, Valore effettivo}
        \end{axis}
    \end{tikzpicture}
    \caption{Percentuale di code coverage dei test implementati}
    \label{fig:10}
\end{figure*}
%--------- FINE GRAFICO -----------%
%\textbf{RTB} \\
% TODO considerazioni finali per RTB \\
%\textbf{PB} \\
% TODO considerazioni finali per PB

\newpage
\subsubsection{25M-BC - Branch coverage}
%--------- GRAFICO -----------%
\begin{figure*}[!h]
    \centering
    \begin{tikzpicture}
        \begin{axis}[
            width  = 0.85*\textwidth,
            height = 8cm,
            ymajorgrids = true,
            symbolic x coords={Sprint 1, Sprint 2, Sprint 3},
            xtick = data,
            ytick = {0, 25, 50, 75, 100},
            ymin=0, ymax=100,
            axis lines*=left,
            legend cell align=left,
            legend style={
                at={(0.5,1.15)},
                anchor=south,
                column sep=0.1ex,
                legend columns=3
            },      
            xticklabel style={rotate=45, anchor=north east, yshift=0ex, xshift=0ex}
            ]
            \addplot[color=opt, style={dashed, line width=3pt}, mark=none] coordinates {(Sprint 1, 100) (Sprint 2, 100) (Sprint 3, 100)};
            \addplot[color=amm, style={dashed, line width=3pt}, mark=none] coordinates {(Sprint 1, 80) (Sprint 2, 80) (Sprint 3, 80)};
            \addplot[color=blue, style={line width=1pt}, mark=none] coordinates {(Sprint 1, 0) (Sprint 2, 0) (Sprint 3, 0)};
            \legend{Valore ottimo, Valore ammissibile, Valore effettivo}
        \end{axis}
    \end{tikzpicture}
    \caption{Percentuale di branch coverage dei test implementati}
    \label{fig:11}
\end{figure*}
%--------- FINE GRAFICO -----------%
%\textbf{RTB} \\
% TODO considerazioni finali per RTB \\
%\textbf{PB} \\
% TODO considerazioni finali per PB

\newpage
\subsubsection{26M-SC - Statement coverage}
%--------- GRAFICO -----------%
\begin{figure*}[!h]
    \centering
    \begin{tikzpicture}
        \begin{axis}[
            width  = 0.85*\textwidth,
            height = 8cm,
            ymajorgrids = true,
            symbolic x coords={Sprint 1, Sprint 2, Sprint 3},
            xtick = data,
            ytick = {0, 25, 50, 75, 100},
            ymin=0, ymax=100,
            axis lines*=left,
            legend cell align=left,
            legend style={
                at={(0.5,1.15)},
                anchor=south,
                column sep=0.1ex,
                legend columns=3
            },      
            xticklabel style={rotate=45, anchor=north east, yshift=0ex, xshift=0ex}
            ]
            \addplot[color=opt, style={dashed, line width=3pt}, mark=none] coordinates {(Sprint 1, 100) (Sprint 2, 100) (Sprint 3, 100)};
            \addplot[color=amm, style={dashed, line width=3pt}, mark=none] coordinates {(Sprint 1, 80) (Sprint 2, 80) (Sprint 3, 80)};
            \addplot[color=blue, style={line width=1pt}, mark=none] coordinates {(Sprint 1, 0) (Sprint 2, 0) (Sprint 3, 0)};
            \legend{Valore ottimo, Valore ammissibile, Valore effettivo}
        \end{axis}
    \end{tikzpicture}
    \caption{Percentuale di statement coverage dei test implementati}
    \label{fig:12}
\end{figure*}
%--------- FINE GRAFICO -----------%
%\textbf{RTB} \\
% TODO considerazioni finali per RTB \\
%\textbf{PB} \\
% TODO considerazioni finali per PB

\newpage
\subsubsection{27M-FD - Failure density}
%--------- GRAFICO -----------%
\begin{figure*}[!h]
    \centering
    \begin{tikzpicture}
        \begin{axis}[
            width  = 0.85*\textwidth,
            height = 8cm,
            ymajorgrids = true,
            symbolic x coords={Sprint 1, Sprint 2, Sprint 3},
            xtick = data,
            ytick = {0, 25, 50, 75, 100},
            ymin=0, ymax=100,
            axis lines*=left,
            legend cell align=left,
            legend style={
                at={(0.5,1.15)},
                anchor=south,
                column sep=0.1ex,
                legend columns=3
            },      
            xticklabel style={rotate=45, anchor=north east, yshift=0ex, xshift=0ex}
            ]
            \addplot[color=opt, style={dashed, line width=3pt}, mark=none] coordinates {(Sprint 1, 100) (Sprint 2, 100) (Sprint 3, 100)};
            \addplot[color=blue, style={line width=1pt}, mark=none] coordinates {(Sprint 1, 0) (Sprint 2, 0) (Sprint 3, 0)};
            \legend{Valore ottimo ed ammissibile, Valore effettivo}
        \end{axis}
    \end{tikzpicture}
    \caption{Percentuale di failure density}
    \label{fig:13}
\end{figure*}
%--------- FINE GRAFICO -----------%
%\textbf{RTB} \\
% TODO considerazioni finali per RTB \\
%\textbf{PB} \\
% TODO considerazioni finali per PB

\newpage
\subsubsection{28M-PTCP - Passed Test Cases Percentage}
%--------- GRAFICO -----------%
\begin{figure*}[!h]
    \centering
    \begin{tikzpicture}
        \begin{axis}[
            width  = 0.85*\textwidth,
            height = 8cm,
            ymajorgrids = true,
            symbolic x coords={Sprint 1, Sprint 2, Sprint 3},
            xtick = data,
            ytick = {0, 25, 50, 75, 100},
            ymin=0, ymax=100,
            axis lines*=left,
            legend cell align=left,
            legend style={
                at={(0.5,1.15)},
                anchor=south,
                column sep=0.1ex,
                legend columns=3
            },      
            xticklabel style={rotate=45, anchor=north east, yshift=0ex, xshift=0ex}
            ]
            \addplot[color=opt, style={dashed, line width=3pt}, mark=none] coordinates {(Sprint 1, 100) (Sprint 2, 100) (Sprint 3, 100)};
            \addplot[color=amm, style={dashed, line width=3pt}, mark=none] coordinates {(Sprint 1, 80) (Sprint 2, 80) (Sprint 3, 80)};
            \addplot[color=blue, style={line width=1pt}, mark=none] coordinates {(Sprint 1, 0) (Sprint 2, 0) (Sprint 3, 0)};
            \legend{Valore ottimo, Valore ammissibile, Valore effettivo}
        \end{axis}
    \end{tikzpicture}
    \caption{Percentuale di casi di test superati}
    \label{fig:14}
\end{figure*}
%--------- FINE GRAFICO -----------%
%\textbf{RTB} \\
% TODO considerazioni finali per RTB \\
%\textbf{PB} \\
% TODO considerazioni finali per PB

\newpage
\subsection{Qualità del processo di Gestione dei rischi}
\subsubsection{29M-NCR - Non-Calculated Risk}
%--------- GRAFICO -----------%
\begin{figure*}[!h]
    \centering
    \begin{tikzpicture}
        \begin{axis}[
            width  = 0.85*\textwidth,
            height = 8cm,
            ymajorgrids = true,
            symbolic x coords={Sprint 1, Sprint 2, Sprint 3},
            xtick = data,
            ymin=0, ymax=5,
            axis lines*=left,
            legend cell align=left,
            legend style={
                at={(0.5,1.15)},
                anchor=south,
                column sep=0.1ex,
                legend columns=3
            },      
            xticklabel style={rotate=45, anchor=north east, yshift=0ex, xshift=0ex}
            ]
            \addplot[color=opt, style={dashed, line width=3pt}, mark=none] coordinates {(Sprint 1, 0) (Sprint 2, 0) (Sprint 3, 0)};
            \addplot[color=amm, style={dashed, line width=3pt}, mark=none] coordinates {(Sprint 1, 3) (Sprint 2, 3) (Sprint 3, 3)};
            \addplot[color=blue, style={line width=1pt}, mark=none] coordinates {(Sprint 1, 1) (Sprint 2, 1) (Sprint 3, 0)};
            \legend{Valore ottimo, Valore ammissibile, Valore effettivo}
        \end{axis}
    \end{tikzpicture}
    \caption{Rischi non calcolati occorsi durante il progetto}
    \label{fig:15}
\end{figure*}
%--------- FINE GRAFICO -----------%
%\textbf{RTB} \\
% TODO considerazioni finali per RTB \\
%\textbf{PB} \\
% TODO considerazioni finali per PB

\newpage
\subsection{Qualità del processo di Gestione della qualità}
\subsubsection{30M-QMS - Quality Metrics Satisfied}
%--------- GRAFICO -----------%
\begin{figure*}[!h]
    \centering
    \begin{tikzpicture}
        \begin{axis}[
            width  = 0.85*\textwidth,
            height = 8cm,
            ymajorgrids = true,
            symbolic x coords={Sprint 1, Sprint 2, Sprint 3},
            xtick = data,
            ytick = {0, 25, 50, 75, 100},
            ymin=0, ymax=100,
            axis lines*=left,
            legend cell align=left,
            legend style={
                at={(0.5,1.15)},
                anchor=south,
                column sep=0.1ex,
                legend columns=3
            },      
            xticklabel style={rotate=45, anchor=north east, yshift=0ex, xshift=0ex}
            ]
            \addplot[color=opt, style={dashed, line width=3pt}, mark=none] coordinates {(Sprint 1, 100) (Sprint 2, 100) (Sprint 3, 100)};
            \addplot[color=amm, style={dashed, line width=3pt}, mark=none] coordinates {(Sprint 1, 90) (Sprint 2, 90) (Sprint 3, 90)};
            \addplot[color=blue, style={line width=1pt}, mark=none] coordinates {(Sprint 1, 43) (Sprint 2, 70) (Sprint 3, 80)};
            \legend{Valore ottimo, Valore ammissibile, Valore effettivo}
        \end{axis}
    \end{tikzpicture}
    \caption{Percentuale di metriche di qualità soddisfatte}
    \label{fig:16}
\end{figure*}
%--------- FINE GRAFICO -----------%
%\textbf{RTB} \\
% TODO considerazioni finali per RTB \\
%\textbf{PB} \\
% TODO considerazioni finali per PB

\newpage
\subsubsection{31M-TE - Efficienza Temporale}
%--------- GRAFICO -----------%
\begin{figure*}[!h]
    \centering
    \begin{tikzpicture}
        \begin{axis}[
            width  = 0.85*\textwidth,
            height = 8cm,
            ymajorgrids = true,
            symbolic x coords={Sprint 1, Sprint 2, Sprint 3},
            xtick = data,
            ymin=0,
            axis lines*=left,
            legend cell align=left,
            legend style={
                at={(0.5,1.15)},
                anchor=south,
                column sep=0.1ex,
                legend columns=3
            },      
            xticklabel style={rotate=45, anchor=north east, yshift=0ex, xshift=0ex}
            ]
            \addplot[color=opt, style={dashed, line width=3pt}, mark=none] coordinates {(Sprint 1, 1) (Sprint 2, 1) (Sprint 3, 1)};
            \addplot[color=amm, style={dashed, line width=3pt}, mark=none] coordinates {(Sprint 1, 3) (Sprint 2, 3) (Sprint 3, 3)};
            \addplot[color=blue, style={line width=1pt}, mark=none] coordinates {(Sprint 1, 2.5) (Sprint 2, 2) (Sprint 3, 4)};
            \legend{Valore ottimo, Valore ammissibile, Valore effettivo}
        \end{axis}
    \end{tikzpicture}
    \caption{Andamento dell'efficienza temporale}
    \label{fig:17}
\end{figure*}
%--------- FINE GRAFICO -----------%
%\textbf{RTB} \\
% TODO considerazioni finali per RTB \\
%\textbf{PB} \\
% TODO considerazioni finali per PB

\newpage
\subsection{Qualità del processo di Pianificazione}
\subsubsection{32M-RSI - Requirements stability index (RSI)}
%--------- GRAFICO -----------%
\begin{figure*}[!h]
    \centering
    \begin{tikzpicture}
        \begin{axis}[
            width  = 0.85*\textwidth,
            height = 8cm,
            ymajorgrids = true,
            symbolic x coords={Sprint 1, Sprint 2, Sprint 3},
            xtick = data,
            ytick = {0, 25, 50, 75, 100},
            ymin=0, ymax=100,
            axis lines*=left,
            legend cell align=left,
            legend style={
                at={(0.5,1.15)},
                anchor=south,
                column sep=0.1ex,
                legend columns=3
            },      
            xticklabel style={rotate=45, anchor=north east, yshift=0ex, xshift=0ex}
            ]
            \addplot[color=opt, style={dashed, line width=3pt}, mark=none] coordinates {(Sprint 1, 100) (Sprint 2, 100) (Sprint 3, 100)};
            \addplot[color=amm, style={dashed, line width=3pt}, mark=none] coordinates {(Sprint 1, 75) (Sprint 2, 75) (Sprint 3, 75)};
            \addplot[color=blue, style={line width=1pt}, mark=none] coordinates {(Sprint 1, 0) (Sprint 2, 80) (Sprint 3, 90)};
            \legend{Valore ottimo, Valore ammissibile, Valore effettivo}
        \end{axis}
    \end{tikzpicture}
    \caption{Percentuale di stabilità dei requisiti}
    \label{fig:18}
\end{figure*}
%--------- FINE GRAFICO -----------%
%\textbf{RTB} \\
% TODO considerazioni finali per RTB \\
%\textbf{PB} \\
% TODO considerazioni finali per PB

\section{Iniziative di automiglioramento per la qualità}
\subsection{Introduzione}
In questa sezione verranno riportate le iniziative di automiglioramento che il nostro gruppo ha deciso di adottare per aumentare la qualità del prodotto e dei processi. Queste iniziative sono state individuate grazie all'esperienza acquisita durante lo svolgimento del progetto e grazie alle valutazioni effettuate sulle attività svolte. \\
Trattandosi per tutti noi della prima esperienza con un progetto di questa portata, è stato necessario un grande numero di tentativi per comprendere al meglio come organizzarci e come svolgere le attività. Questo ci ha permesso di capire quali sono stati i punti di forza e i punti deboli del nostro lavoro e di individuare le aree in cui è possibile migliorare. \\
Per ciascuna delle difficoltà riscontrate verranno indicate:
\begin{itemize}
    \item fase del progetto in cui si è verificato il problema;
    \item descrizione del problema;
    \item contromisura adottata per risolvere il problema evidenziato.
\end{itemize}

\subsection{Problemi rilevati ed iniziative adottate}
\begin{itemize}
    \item \textbf{Organizzazione delle riunioni}
    \begin{itemize}
        \item \textbf{Fase del progetto}: iniziale;
        \item \textbf{Descrizione}: nelle prime settimane di lavoro, a partire dalla formazione dei gruppi sino ai primi Diari di bordo, si è riscontrata una certa difficoltà nell'organizzazione delle riunioni causata dai vari impegni di ciascun membro (lezioni diverse in orari diversi, lavoro per alcuni, impegni personali) e soprattutto alimentata dalle diverse riunioni che si accumulavano (SAL con l'azienda prima e Diari di bordo poi) portando a una certa confusione e a un rallentamento delle attività;
        \item \textbf{Contromisura}: abbiamo deciso di effettuare le riunioni a distanza tramite la piattaforma \textit{Discord} e di fissare un giorno e un orario durante la settimana per ciascuna tipologia di incontro in maniera tale da rispettare le disponibilità di ogni membro; qualora qualcuno, per impegni di natura eccezionale, non abbia modo di essere presente potrà successivamente informarsi sui contenuti trattati attraverso i verbali che verranno redatti e messi a disposizione di tutti.
    \end{itemize}
    \item \textbf{Suddivisione compiti}
    \begin{itemize}
        \item \textbf{Fase del progetto}: iniziale;
        \item \textbf{Descrizione}: all'inizio del progetto si è riscontrata una certa difficoltà nella suddivisione dei compiti a causa della mancanza di esperienza e della poca conoscenza delle competenze possedute da ciascuno. È risultato dunque difficile il bilanciamento delle mansioni e si sono verificati più volte casi in cui alcuni membri sono stati in grado di completare le attività a loro assegnate in anticipo, e casi opposti in cui il lavoro da svolgere è risultato eccessivo e difficilmente completabile entro i tempi prestabiliti;
        \item \textbf{Contromisura}: abbiamo quindi deciso, come suggerito anche dal professor Vardanega al primo Diario di bordo, di non assegnare preventivamente tutti i compiti da svolgere a ciascun membro, ma piuttosto di metterli in un contenitore condiviso (abbiamo deciso di usare le annotazioni di \textit{ClickUp}) e di permettere a ciascun membro di prendere in autonomia i compiti da svolgere, così che chiunque finisca in anticipo possa prenderne altri; in questo modo siamo riusciti a svolgere le attività in modo più equo e a completare i compiti entro i tempi prestabiliti.
    \end{itemize}
    \item \textbf{Familiarità con le tecnologie}
    \begin{itemize}
        \item \textbf{Fase del progetto}: intermedia;
        \item \textbf{Descrizione}: durante lo svolgimento del progetto ci siamo resi conto che la mancanza di familiarità con le tecnologie utilizzate (in particolare con \textit{Docker}, \textit{Grafana} e \textit{Clickhouse}) ha rallentato inizialmente l'attività di sviluppo e ha portato a un aumento del carico di lavoro per alcuni membri del gruppo;
        \item \textbf{Contromisura}: abbiamo deciso di organizzare un incontro di formazione in cui i membri più esperti hanno spiegato ai meno esperti il funzionamento di \textit{Docker} e le modalità di utilizzo. Inoltre, abbiamo deciso di utilizzare la funzionalità di \textit{pair programming} per permettere ai membri meno esperti di lavorare a stretto contatto con quelli più esperti e di apprendere da loro.
    \end{itemize}
\end{itemize}

\subsection{Considerazioni finali}
Fin da subito il nostro gruppo si è posto come obiettivo principale quello di dotarsi di un \textit{Way of Working} preciso e ben definito, di pianificare ogni singola attività e di prevedere tutte le possibili difficoltà incontrabili durante lo svolgimento del progetto. Questo per cercare di prevenire i problemi e di fornire delle contromisure efficaci per affrontarli. \\
Inizialmente si sono presentate delle difficoltà dovute all'inesperienza del gruppo in ambito organizzativo. Tuttavia, grazie alla familiarizzazione ottenuta tramite lo svolgimento del progetto e grazie ai consigli e suggerimenti che ci sono stati forniti dai professori e dall'azienda proponente, siamo riusciti a individuare i problemi e a mettere in atto delle contromisure per risolverli. \\
Questo ci ha permesso di migliorare notevolmente la qualità del nostro lavoro e di svolgere le attività in modo più efficiente e più equo. Nonostante ciò siamo anche consapevoli che ci sono ancora molti aspetti su cui possiamo progredire e che ci sono ancora molte iniziative di automiglioramento che possiamo adottare. Siamo convinti che, se continueremo a lavorare con lo stesso impegno e la stessa determinazione che abbiamo dimostrato finora, saremo in grado di ottenere risultati di qualità superiore.
\end{document}